% Generated mechanically from the frozen introduction.tex target.
% Do not edit this generated file; edit the bound locale source instead.

% OK, start here.
%

\setcounter{chapter}{0}
\chapter{引言}



\phantomsection
\label{introduction-section-phantom}

\begin{verbatim}
Copyright (C) 2005 -- 2025 Johan de Jong
Permission is granted to copy, distribute and/or modify this
document under the terms of the GNU Free Documentation License,
Version 1.2 or any later version published by the Free Software
Foundation; with no Invariant Sections, no Front-Cover Texts,
and no Back-Cover Texts. A copy of the license is included in
the section entitled "GNU Free Documentation License".
\end{verbatim}



\section{概述}
\label{introduction-section-overview}

\noindent
除 Laumon 与 Moret-Bailly 的著作（见 \cite{LM-B}）以及 Fulton 等人尚在进行的工作之外，
我们认为，仍有必要编写一部关于代数栈及定义代数栈所需代数几何的开源教材。
Stacks Project 试图从交换代数出发奠定基础，继而经过概形与代数空间理论，
最终为代数栈理论建立一套全面的基础。

\medskip\noindent
我们设想本材料主要供在线阅读，因为它的一项关键特色是超链接：
这些链接能够迅速跳转至散布在许多不同页面上的内部引用。
如果浏览器使用嵌入式 PDF 或 DVI 查看器，跨文件链接应当可以正常工作。

\medskip\noindent
本项目由多人协作完成，我们诚邀读者参与。
阅读时如发现任何排印错误或内容错误，或有建议、补充材料或例子，
请发送电子邮件至
\href{mailto:stacks.project@gmail.com}{stacks.project@gmail.com}。
你可以下载包含全部源文件的 tar 归档包，解压后运行 make，
并在本地使用 DVI 或 PDF 查看器阅读。
也欢迎直接编辑这些 LaTeX 文件，并将改进稿通过电子邮件发给我们。


\section{贡献说明}
\label{introduction-section-attribution}

\noindent
本书涵盖范围极广；要准确而简洁地注明所有促成我们在此试图阐述之理论发展的数学家，
是一项艰巨的任务。我们希望最终能够引起足够的社群兴趣，找到愿意为每一章撰写历史评注
部分的贡献者。

\medskip\noindent
对本项目作出贡献者的姓名列于本书版本的扉页，并另有
\href{https://stacks.math.columbia.edu/tex/CONTRIBUTORS}{在线名单}。
这里列举若干主要贡献：
\begin{enumerate}
\item Jarod Alper 贡献了一章讨论代数栈相关文献的内容，见
《文献指南》第 \ref{guide-section-short-introductions} 节。
\item Bhargav Bhatt 撰写了概形的平展态射一章的初稿，见
《平展态射》第 \ref{etale-section-introduction} 节。
\item Bhargav Bhatt 撰写了
《代数进阶》第 \ref{more-algebra-section-formal-glueing} 节的初稿。
\item Kiran Kedlaya 贡献了
《下降》第 \ref{descent-section-descent-universally-injective} 节的初稿。
\item 下列内容的初稿：
\begin{enumerate}
\item 《代数》第 \ref{algebra-section-oka-families} 节，
\item 《内射对象》第 \ref{injectives-section-baer} 节，以及
\item 域这一章，见
《域》第 \ref{fields-section-introduction} 节，
\end{enumerate}
均取自
\href{https://math.uchicago.edu/~amathew/cr.html}{The CRing Project}，
由 Akhil Mathew 等人惠允使用。
\item Alex Perry 撰写了关于投射模和 Mittag--Leffler 模的材料，
其中包括
《代数》定理 \ref{algebra-theorem-ffdescent-projectivity} 的证明。
\item Alex Perry 撰写了 Schlessinger 与 Rim 意义下的形变理论一章，见
《形式形变理论》第 \ref{formal-defos-section-introduction} 节。
\item Thibaut Pugin、Zachary Maddock 与 Min Lee 为一门课程记录了讲义，
这些讲义构成了平展上同调一章与迹公式一章的基础。见
《平展上同调》第 \ref{etale-cohomology-section-introduction} 节及
《迹公式》第 \ref{trace-section-introduction} 节。
\item David Rydh 提出了许多有益意见，指出了若干错误，
在剩余 gerbe 相关材料上提供了关键帮助，并且是以下材料的最初提出者：
《空间中的群胚进阶》第
\ref{spaces-more-groupoids-section-finite} 节和
\ref{spaces-more-groupoids-section-etale-localize} 节。
\item Burt Totaro 贡献了《例》第
\ref{examples-section-non-descending-property-projective} 节、
\ref{examples-section-non-effective-descent-projective} 节，以及
《栈的性质》第 \ref{stacks-properties-section-dimension} 节。
\item pro-平展上同调一章（见
《pro-平展上同调》第 \ref{proetale-section-introduction} 节）
取自 Bhargav Bhatt 与 Peter Scholze 的一篇论文。
\item Bhargav Bhatt 贡献了《例》第
\ref{examples-section-non-algebraic-hom-stack} 节和
\ref{examples-section-flat-not-colimit-flat-finitely-presented} 节。
\item Ofer Gabber 发现了错误并提供了修正，还贡献了
《簇》引理 \ref{varieties-lemma-image-connected-component}、
《形式空间》引理 \ref{formal-spaces-lemma-completion-countably-indexed}、
《群胚进阶》第 \ref{more-groupoids-section-ind-quasi-affine} 节中的材料、
《空间的性质》第 \ref{spaces-properties-section-fpqc} 节的主要结果，
以及《平坦性进阶》命题
\ref{flat-proposition-finite-type-injective-into-flat-mod-m} 的证明。
\item J\'anos Koll\'ar 贡献了
《代数》引理 \ref{algebra-lemma-hart-serre-loc-thm} 和
《局部上同调》命题 \ref{local-cohomology-proposition-kollar}。
\item Kiran Kedlaya 撰写了
《代数进阶》第 \ref{more-algebra-section-beauville-laszlo} 节的初稿。
\item Matthew Emerton、Toby Gee 与 Brandon Levin 贡献了若干关于加厚的结果，
尤其包括《栈的态射进阶》中的引理
\ref{stacks-more-morphisms-lemma-reduced-diagonal}、
\ref{stacks-more-morphisms-lemma-thickening-diagonals} 和
\ref{stacks-more-morphisms-lemma-thickening-properties}。
\item Lena Min Ji 撰写了
《代数进阶》第 \ref{more-algebra-section-principal-radical-ideals} 节的初稿。
\item Matthew Emerton 与 Toby Gee 撰写了《栈的几何》第
\ref{stacks-geometry-section-multiplicities} 节和
\ref{stacks-geometry-section-dimension-of-algebraic-stacks} 节的初稿。
\end{enumerate}
